\documentclass{subfiles}

\begin{document}

\subsection{Bermain dengan Matriks}
Karena terlalu fokus dengan hari raya, Holil hampir lupa mengerjakan tugasnya. Holil memiliki tugas materi matriks. Tugas Holil adalah melakukan pencerminan dan rotasi. Pertama, Holil harus melakukan rotasi 90 derajat pada matriks. Setelah itu, dilakukan pencerminan terhadap sumbu X.

\inputfigure{res/m3i1.png}{0.21}{Keluaran {\ioI} Genap Program Bermain dengan Matriks}{fig:m3:i1}

\inputfigure{res/m3i2.png}{0.21}{Keluaran {\ioI} Ganjil Program Bermain dengan Matriks}{fig:m3:i2}

\subsection{Enkripsi dan Dekripsi}

Setelah mengerjakan tugas, Holil lupa bahwa dia belum mengucapkan minal aidin walfaidzin kepada teman-temannya. Namun, Holil ingin ucapan itu ia sampaikan secara rahasia, agar temannya mengeluarkan usaha lebih dalam memahami pesannya. Holil pun memutuskan untuk menggunakan 2 metode kriptografi yaitu \textit{Reverse cipher} dan Caesar {\cipher}.

Holil melakukan enkripsi pada pesannya dengan menggunakan {\key} ``5'' bergeser ke kanan untuk melakukan Caesar {\cipher}.

\inputfigure{res/m3i3.png}{0.3}{Hasil Penjalanan Program Enkripsi dan Dekripsi}{fig:m3:i3}
\newpage
\subsectionf{\textit{Blast the Fireworks!}}

Holil masih memiliki stok kembang api. Kembang api yang dimiliki Holil memiliki kekuatan ledakan yang berbeda-beda. Ia pun ingin menghabiskan stok kembang apinya. Pertama-tama, Holil akan menentukan jumlah kembang api yang ingin diledakkan, kemudian kembang api akan memiliki \textit{range} ledakan antara 0-9. Setelah itu, Holil menentukan satu kembang api yang akan ia ledakkan terlebih dahulu.

\inputfigure{res/m3i4.png}{0.3}{Hasil Penjalanan Program \textit{Blast the Fireworks}}{fig:m3:i4}

\end{document}
